\selectlanguage{russian}
\section{Эксперименты}
\label{sec:Experiments} \index{Experiments}

Эксперименты предназначены для проверки двух аспектов предложенного подхода: качества
восстановления скрытого сопоставления и качества оценки условной плотности. Исследуется
зависимость результатов от уровня шума $\epsilon$, размера группы $M$ и числа групп $N$.
В качестве внешней точки сравнения используется энтропийный транспортный оператор Байера
и др.~\cite{beier2024transfer}, решающий ту же статистическую задачу.

\subsection{Синтетические данные}
\label{ssec:exp-data}

Используется общий с референсной работой синтетический пример на одномерном торе
$\mathbb{S}^1=[0,1)$. Величина $x$ имеет равномерное распределение, а условный закон
$Y\mid X=x$ представляет собой бимодальную обёрнутую гауссову смесь:
\begin{equation}
p^*(x)=\mathcal{U}[0,1),
\qquad
p^*(y\mid x)=\tfrac12\tilde{\mathcal{N}}(x,\epsilon^2)
+\tfrac12\tilde{\mathcal{N}}(x+0.3,\epsilon^2).
\label{eq:exp-generator}
\end{equation}
Здесь $\epsilon$ обозначает стандартное отклонение шума. Для каждой группы независимо
генерируются $M$ пар, после чего значения $y$ переупорядочиваются случайной перестановкой.
Алгоритм получает только два набора значений и не использует истинное соответствие при
обучении. Набор данных содержит $N$ таких групп.

\begin{figure}[H]
\centering
\includegraphics[width=0.85\linewidth]{data_orig_generation_torus.png}
\caption{Синтетические данные на торе при $\epsilon=0.05$. Слева показаны $10^3$ пар
$(x,y)$, справа --- их двумерная гистограмма. Полосы соответствуют компонентам
$y\approx x$ и $y\approx x+0.3$ с периодическим продолжением.}
\label{fig:data-torus}
\end{figure}

\subsection{Метрики и протокол}
\label{ssec:exp-protocol}

\paragraph{Ошибка восстановления плотности.}
Поскольку плотность равномерного распределения $X$ на $[0,1)$ равна единице, на
рассматриваемой области выполняется $p(x,y)=p(y\mid x)$. Истинная и восстановленная
плотности сравниваются на равномерной сетке размера $E\times E$:
\begin{equation}
L_2=\frac{1}{E}\left\|p^*(x,y)-\widehat p_\theta(x,y)\right\|_2
=\left(
\frac{1}{E^2}\sum_{i=1}^{E}\sum_{j=1}^{E}
\bigl(p^*(x_i,y_j)-\widehat p_\theta(x_i,y_j)\bigr)^2
\right)^{1/2}.
\label{eq:l2-metric}
\end{equation}
Меньшее значение соответствует более точному восстановлению. Для нейросетевой модели
плотность определяется условной смесью~\eqref{eq:decoder-mixture}; для глобальной смеси
сдвигов используется
$\widehat p_\theta(y\mid x)=\sum_c\pi_c
\tilde{\mathcal{N}}(y\mid x+s_c,\epsilon^2)$.

\paragraph{Точность локальных назначений.}
Для диагностического эксперимента истинная перестановка сохраняется только на этапе оценки.
Предсказанный индекс определяется построчным $\arg\max$ мягкой матрицы назначений, после
чего вычисляется доля совпавших индексов. Эта метрика оценивает локальные назначения и, в
отличие от Hungarian matching, сама по себе не гарантирует биективность всей предсказанной
перестановки.

\paragraph{Протокол.}
Каждая точка кривых соответствует медиане по пяти независимым запускам. Полупрозрачная
область показывает межквартильный размах. Основные гиперпараметры приведены в
Таблице~\ref{tab:exp-params}. Нейросетевая модель обучается стохастическим градиентным
методом, а глобальная смесь сдвигов --- полнобатчевым модифицированным EM с
Sinkhorn-аппроксимацией E-шага и аналитическим M-шагом. Это два подхода с разными
параметризациями: модель сдвигов использует заранее известную структуру синтетического
генератора, тогда как CatVAE задаёт более общий условный закон нейронными сетями.

\begin{table}[H]
\centering
\caption{Основные параметры экспериментов.}
\label{tab:exp-params}
\begin{tabular}{lll}
\toprule
Параметр & CatVAE & Глобальная смесь сдвигов \\
\midrule
Уровни шума $\epsilon$ & \multicolumn{2}{l}{$\{0.01,0.025,0.05\}$} \\
Число компонент $K$ & \multicolumn{2}{l}{$2$} \\
Сетка метрики $E$ & \multicolumn{2}{l}{$100$} \\
Число запусков & \multicolumn{2}{l}{$5$ (медиана и IQR)} \\
Число итераций & $5000$ ($M,N$); $2500$ ($\epsilon$) & $200$ \\
Температура $\tau$ & \multicolumn{2}{l}{$1$ (фиксирована)} \\
Число Sinkhorn-итераций & $50$ & $50$ \\
Шаг обучения & $10^{-3}$ & --- \\
Число групп в опыте по $M$ & $20$ & $20$ \\
Исследованный диапазон $M$ & \multicolumn{2}{l}{$1\dots512$} \\
Диапазон $N$ при $M=20$ & \multicolumn{2}{l}{$1\dots407$} \\
\bottomrule
\end{tabular}
\end{table}

Референсные значения взяты из результатов Байера и др. для того же синтетического
распределения и той же $L_2$-метрики. Поэтому сравнение характеризует качество оценок на
общей задаче, но не является изолированным сравнением оптимизаторов или времени работы.

\subsection{Восстановление скрытого сопоставления}
\label{ssec:exp-recovery}

На Рисунке~\ref{fig:recovery} показаны результаты CatVAE при $M=2$. На левых панелях
значения соединены по наблюдаемому после перемешивания индексу. На правых панелях пары
образованы по построчному $\arg\max$ мягкой матрицы назначений. Зелёный цвет соответствует
совпадению с сохранённым истинным индексом, красный --- ошибке.

\begin{figure}[H]
\centering
\includegraphics[width=0.82\linewidth]{sigma_0.01_sgd_spread_recovery.png}\\[2pt]
\includegraphics[width=0.82\linewidth]{sigma_0.025_sgd_spread_recovery.png}\\[2pt]
\includegraphics[width=0.82\linewidth]{sigma_0.05_sgd_spread_recovery.png}
\caption{Локальные назначения CatVAE при $M=2$ для $\epsilon=0.01$ (сверху),
$\epsilon=0.025$ (в центре) и $\epsilon=0.05$ (снизу). Слева показано исходное
переупорядочивание, справа --- назначения по построчному $\arg\max$.}
\label{fig:recovery}
\end{figure}

Численные значения приведены в Таблице~\ref{tab:recovery-acc}. При случайном выборе одного
из двух индексов ожидаемая точность составляет $50\%$. Полученные значения существенно выше
этого уровня. С ростом $\epsilon$ компоненты условной смеси сильнее перекрываются, а
наблюдаемая точность назначений снижается. Этот эксперимент подтверждает, что матрица
назначений содержит информацию об исходных парах, но не доказывает однозначную
восстановимость каждой перестановки при произвольных $M$ и $\epsilon$.

\begin{table}[H]
\centering
\caption{Точность локальных назначений CatVAE при $M=2$.}
\label{tab:recovery-acc}
\begin{tabular}{cc}
\toprule
Уровень шума $\epsilon$ & Точность \\
\midrule
$0.01$ & $98.0\%$ \\
$0.025$ & $95.6\%$ \\
$0.05$ & $89.7\%$ \\
\bottomrule
\end{tabular}
\end{table}

\subsection{Влияние размера группы \texorpdfstring{$M$}{M}}
\label{ssec:exp-M}

На Рисунке~\ref{fig:varying-M} показана зависимость $L_2$-ошибки от числа точек в группе
при фиксированных $N=20$. Сплошные линии соответствуют двум исследуемым параметризациям,
штриховые --- референсному методу. Цвет обозначает уровень порождающего шума.

\begin{figure}[H]
\centering
\begin{minipage}{0.49\textwidth}
\centering
\includegraphics[width=\linewidth]{varying_M_sgd_spread.png}\\
{\footnotesize (a) CatVAE}
\end{minipage}
\hfill
\begin{minipage}{0.49\textwidth}
\centering
\includegraphics[width=\linewidth]{varying_M_em_tau1.0.png}\\
{\footnotesize (b) Глобальная смесь сдвигов}
\end{minipage}
\caption{$L_2$-ошибка в зависимости от $M$ при $N=20$. Сплошные кривые показывают
медиану по пяти запускам, области --- IQR, штриховые кривые --- референс. Синий,
оранжевый и зелёный цвета соответствуют $\epsilon=0.01$, $0.025$ и $0.05$.}
\label{fig:varying-M}
\end{figure}

При малых и умеренных $M$ глобальная смесь сдвигов даёт меньшую ошибку, чем референсный
метод. Этот результат ожидаем для параметризации, совпадающей по форме с генератором данных:
требуется оценить только два общих сдвига и их веса. Он не позволяет приписать преимущество
самому модифицированному EM независимо от выбранной модели.

При $\tau=1$ ошибка смеси сдвигов не возрастает с увеличением $M$: во всём исследованном
диапазоне она остаётся существенно ниже референсных значений. Следовательно, сам по себе
рост матрицы назначений в данном эксперименте не приводит к наблюдавшемуся ранее излому.
Зависимость этого эффекта от температуры отдельно исследуется ниже.

Ошибка CatVAE остаётся сопоставимой с референсом, но в большинстве режимов превышает ошибку
специализированной смеси сдвигов. Возможной причиной является большая гибкость условной
модели и, следовательно, более сложная совместная оценка плотности и назначений.

\subsection{Влияние числа групп \texorpdfstring{$N$}{N}}
\label{ssec:exp-N}

На Рисунке~\ref{fig:varying-N} показана зависимость ошибки от числа групп при $M=20$.
Сплошные кривые соответствуют трём уровням шума исследуемых моделей. Штриховые кривые
показывают референс при $\epsilon=0.05$ и трёх значениях силы энтропийной регуляризации
$\varepsilon$.

\begin{figure}[H]
\centering
\begin{minipage}{0.49\textwidth}
\centering
\includegraphics[width=\linewidth]{varying_N_sgd_spread.png}\\
{\footnotesize (a) CatVAE}
\end{minipage}
\hfill
\begin{minipage}{0.49\textwidth}
\centering
\includegraphics[width=\linewidth]{varying_N_em_tau1.0.png}\\
{\footnotesize (b) Глобальная смесь сдвигов}
\end{minipage}
\caption{$L_2$-ошибка в зависимости от $N$ при $M=20$. Сплошные кривые соответствуют
$\epsilon\in\{0.01,0.025,0.05\}$; штриховые кривые --- референсу при $\epsilon=0.05$ и
$\varepsilon\in\{0.001,0.0025,0.01\}$. Показаны медиана и IQR по пяти запускам.}
\label{fig:varying-N}
\end{figure}

При $\tau=1$ ошибка глобальной смеси сдвигов уменьшается с ростом $N$ и становится близкой
к нулю при всех трёх уровнях шума. Это соответствует низкой размерности правильно
специфицированной модели: дополнительные группы уточняют оценки общих сдвигов и весов.
В частности, плато при $\epsilon=0.05$, возникающее при меньших температурах, не является
неизбежным следствием перекрытия компонент.

При малых $N$ отдельные кривые немонотонны, а межквартильный размах увеличивается. Это
указывает на высокую чувствительность оценки к конкретной выборке и инициализации при
дефиците данных. Гарантии монотонности классического EM здесь неприменимы по двум причинам:
E-шаг модифицирован Sinkhorn-аппроксимацией, и сравниваются разные обучающие выборки, а не
соседние итерации одного запуска. Поэтому немонотонность по $N$ рассматривается как
эмпирическая нестабильность малой выборки, а не как следствие свойства ELBO.

В совокупности эксперименты показывают, что обе параметризации способны восстанавливать
условную плотность из непарных групп. Специализированная смесь сдвигов наиболее точна при
совпадении её структуры с генератором, тогда как CatVAE не использует предположение о
глобально постоянных сдвигах. Для модифицированного EM при низкой температуре качество
ограничивается устойчивостью Sinkhorn-аппроксимации E-шага, как показывает следующий опыт.

\paragraph{Влияние температуры.}
Для модифицированного EM дополнительно проведены эксперименты при
$\tau\in\{0.05,0.1,0.2,0.5\}$. Все остальные параметры совпадают с основным протоколом.
Результаты для каждого значения $\tau$ образуют отдельный столбец на
Рисунке~\ref{fig:em-temperature}: верхний ряд соответствует варьированию $M$ при $N=20$,
нижний --- варьированию $N$ при $M=20$.

\begin{figure}[H]
\centering
\begin{tabular}{@{}cccc@{}}
\includegraphics[width=0.24\textwidth]{varying_M_em_tau0.05.png} &
\includegraphics[width=0.24\textwidth]{varying_M_em_tau0.1.png} &
\includegraphics[width=0.24\textwidth]{varying_M_em_tau0.2.png} &
\includegraphics[width=0.24\textwidth]{varying_M_em_tau0.5.png} \\
{\footnotesize $\tau=0.05$} & {\footnotesize $\tau=0.1$} &
{\footnotesize $\tau=0.2$} & {\footnotesize $\tau=0.5$} \\[2pt]
\includegraphics[width=0.24\textwidth]{varying_N_em_tau0.05.png} &
\includegraphics[width=0.24\textwidth]{varying_N_em_tau0.1.png} &
\includegraphics[width=0.24\textwidth]{varying_N_em_tau0.2.png} &
\includegraphics[width=0.24\textwidth]{varying_N_em_tau0.5.png}
\end{tabular}
\caption{Чувствительность глобальной смеси сдвигов к температуре $\tau$. Верхний ряд:
$L_2$-ошибка в зависимости от $M$ при $N=20$; нижний ряд: ошибка в зависимости от $N$
при $M=20$. Один столбец соответствует одному значению $\tau$. Обозначения кривых
совпадают с Рисунками~\ref{fig:varying-M} и~\ref{fig:varying-N}.}
\label{fig:em-temperature}
\end{figure}

В эксперименте по $M$ при всех исследованных $\tau<1$ ошибка после некоторого порогового
размера группы возрастает к одним и тем же уровням; при малых $\tau$ хорошо различимо
плато. Это согласуется с формой
E-шага~\eqref{eq:em-e}: деление $L$ на малую температуру усиливает различия попарных
логарифмических правдоподобий, поэтому план $q$ становится почти дискретным. Случайное
небольшое преимущество одной компоненты тогда даёт ей почти всю ответственность
$r_{mkc}=q_{mk}w_{mkc}$. Обновления~\eqref{eq:em-m-pi} и~\eqref{eq:em-m-shift} усиливают
это преимущество, и вторая компонента теряет вес либо совмещается с первой. В результате
оценка условной плотности фактически содержит только одну из двух истинных мод. При меньшей
$\tau$ эта положительная обратная связь включается резче, что объясняет более резкий излом
кривых. Совпадение наблюдаемых уровней плато с однокомпонентным пределом служит
дополнительным свидетельством такого коллапса.

Этот предел можно вычислить явно. Обозначим плотность обёрнутого нормального распределения
с нулевым средним через
$g_\epsilon(t)=\sum_{k\in\mathbb{Z}}\mathcal{N}(t+k\mid0,\epsilon^2)$ и положим
$d=0.3$. Истинная плотность разности $y-x$ равна
$\tfrac12 g_\epsilon(t)+\tfrac12 g_\epsilon(t-d)$. Если оценка сохраняет только одну моду,
например $\widehat p(t)=g_\epsilon(t)$, то
\begin{align}
L_{2,\mathrm{coll}}^2
&=\frac14\int_0^1\bigl(g_\epsilon(t-d)-g_\epsilon(t)\bigr)^2\,dt \\
&=\frac12\sum_{\ell\in\mathbb{Z}}
e^{-4\pi^2\epsilon^2\ell^2}\bigl(1-\cos(2\pi\ell d)\bigr).
\label{eq:collapsed-mixture-l2}
\end{align}
Поскольку $d$ велик по сравнению с исследуемыми $\epsilon$, перекрытие мод мало и
\begin{equation}
L_{2,\mathrm{coll}}
\simeq \frac{1}{2\pi^{1/4}\sqrt{\epsilon}}.
\label{eq:collapsed-mixture-l2-approx}
\end{equation}
Отсюда для дисперсий $\epsilon^2\in\{10^{-4},6.25\cdot10^{-4},2.5\cdot10^{-3}\}$
получаются значения $L_{2,\mathrm{coll}}\approx3.76$, $2.38$ и $1.68$ соответственно.
Именно около этих уровней располагаются синяя, оранжевая и зелёная кривые после перехода.
